% GENERATED FILE. Edit canonical lessons, then run npm run generate:books.
% canonical-source-hash: fnv1a64:62f78eb2ebeedcf2
% canonical-lessons: SA-C251-prthivi, SA-C251-jagat, SA-C251-lokah, SA-C251-prakrtih, SA-C251-pankah

\chapter{Earth, Universe, World, Nature, Mud}
\label{ch:sa-earth-universe-world-nature-mud}
\hlchaptermodality{251}

\begin{chapteropening}
\textbf{By the end of this chapter:} \emph{I can say the earth, the universe, the world, nature and mud.}

Complete the last lesson of chapter 251: I can say the earth, the universe, the world, nature and mud.
\end{chapteropening}

\section[\texorpdfstring{\sk{पृथिवी} (pṛthivī)}{पृथिवी (pṛthivī)}]{\sk{पृथिवी} (pṛthivī) — the earth}

\label{lesson:SA-C251-prthivi}



\begin{quote}
Before the new one: say the Sanskrit for a building, then the Sanskrit for lodging.
\end{quote}

\subsection*{You'll want to know: \sk{पृथिवी}}
\textbf{\sk{पृथिवी}} — \emph{pṛthivī} — \textquotedblleft{}the earth\textquotedblright{}.

\begin{grammarlens}[title={What you've built}]
One more word for this chapter.
\end{grammarlens}

\subsection*{Your turn}
\begin{itemize}
\raggedright
  \item \emph{Say it:} \emph{pṛthivī}
  \item \emph{Say it:} \emph{pṛthivī}, once more
  \item \emph{Say it:} say \emph{pṛthivī}
  \item \emph{Recall:} say the Sanskrit for a building, then the Sanskrit for lodging, then say \emph{pṛthivī} again
\end{itemize}

\begin{tcolorbox}[breakable,colback=teal!4,colframe=teal!35!black,title={Before you move on}]
What does \sk{पृथिवी} mean? (\textquotedblleft{}The earth\textquotedblright{}.) Say it once more.
\end{tcolorbox}
\section[\texorpdfstring{\sk{जगत्} (jagat)}{जगत् (jagat)}]{\sk{जगत्} (jagat) — the universe, the world}

\label{lesson:SA-C251-jagat}



\begin{quote}
Before the new one: say the Sanskrit for lodging, then the Sanskrit for the earth.
\end{quote}

\subsection*{You'll want to know: \sk{जगत्}}
\textbf{\sk{जगत्}} — \emph{jagat} — \textquotedblleft{}the universe, the world\textquotedblright{}.

\begin{grammarlens}[title={What you've built}]
One more word for this chapter.
\end{grammarlens}

\subsection*{Your turn}
\begin{itemize}
\raggedright
  \item \emph{Say it:} \emph{jagat}
  \item \emph{Say it:} \emph{jagat}, once more
  \item \emph{Say it:} say \emph{jagat}
  \item \emph{Recall:} say the Sanskrit for lodging, then the Sanskrit for the earth, then say \emph{jagat} again
\end{itemize}

\begin{tcolorbox}[breakable,colback=teal!4,colframe=teal!35!black,title={Before you move on}]
What does \sk{जगत्} mean? (\textquotedblleft{}The universe, the world\textquotedblright{}.) Say it once more.
\end{tcolorbox}
\section[\texorpdfstring{\sk{लोकः} (lokaḥ)}{लोकः (lokaḥ)}]{\sk{लोकः} (lokaḥ) — the world, people}

\label{lesson:SA-C251-lokah}



\begin{quote}
Before the new one: say the Sanskrit for the earth, then the Sanskrit for the universe.
\end{quote}

\subsection*{You'll want to know: \sk{लोकः}}
\textbf{\sk{लोकः}} — \emph{lokaḥ} — \textquotedblleft{}the world, people\textquotedblright{}.

\begin{grammarlens}[title={What you've built}]
One more word for this chapter.
\end{grammarlens}

\subsection*{Your turn}
\begin{itemize}
\raggedright
  \item \emph{Say it:} \emph{lokaḥ}
  \item \emph{Say it:} \emph{lokaḥ}, once more
  \item \emph{Say it:} say \emph{lokaḥ}
  \item \emph{Recall:} say the Sanskrit for the earth, then the Sanskrit for the universe, then say \emph{lokaḥ} again
\end{itemize}

\begin{tcolorbox}[breakable,colback=teal!4,colframe=teal!35!black,title={Before you move on}]
What does \sk{लोकः} mean? (\textquotedblleft{}The world, people\textquotedblright{}.) Say it once more.
\end{tcolorbox}
\section[\texorpdfstring{\sk{प्रकृतिः} (prakṛtiḥ)}{प्रकृतिः (prakṛtiḥ)}]{\sk{प्रकृतिः} (prakṛtiḥ) — nature}

\label{lesson:SA-C251-prakrtih}



\begin{quote}
Before the new one: say the Sanskrit for the universe, then the Sanskrit for the world.
\end{quote}

\subsection*{You'll want to know: \sk{प्रकृतिः}}
\textbf{\sk{प्रकृतिः}} — \emph{prakṛtiḥ} — \textquotedblleft{}nature\textquotedblright{}.

\begin{grammarlens}[title={What you've built}]
One more word for this chapter.
\end{grammarlens}

\subsection*{Your turn}
\begin{itemize}
\raggedright
  \item \emph{Say it:} \emph{prakṛtiḥ}
  \item \emph{Say it:} \emph{prakṛtiḥ}, once more
  \item \emph{Say it:} say \emph{prakṛtiḥ}
  \item \emph{Recall:} say the Sanskrit for the universe, then the Sanskrit for the world, then say \emph{prakṛtiḥ} again
\end{itemize}

\begin{tcolorbox}[breakable,colback=teal!4,colframe=teal!35!black,title={Before you move on}]
What does \sk{प्रकृतिः} mean? (\textquotedblleft{}Nature\textquotedblright{}.) Say it once more.
\end{tcolorbox}
\section[\texorpdfstring{\sk{पंकः} (paṅkaḥ)}{पंकः (paṅkaḥ)}]{\sk{पंकः} (paṅkaḥ) — mud}

\label{lesson:SA-C251-pankah}



\begin{quote}
Before the new one: say the Sanskrit for the world, then the Sanskrit for nature.
\end{quote}

\subsection*{You'll want to know: \sk{पंकः}}
\textbf{\sk{पंकः}} — \emph{paṅkaḥ} — \textquotedblleft{}mud\textquotedblright{}.

\begin{grammarlens}[title={What you've built}]
One more word for this chapter.
\end{grammarlens}

\subsection*{Your turn}
\begin{itemize}
\raggedright
  \item \emph{Say it:} \emph{paṅkaḥ}
  \item \emph{Say it:} \emph{paṅkaḥ}, once more
  \item \emph{Say it:} say \emph{paṅkaḥ}
  \item \emph{Recall:} say the Sanskrit for the world, then the Sanskrit for nature, then say \emph{paṅkaḥ} again
\end{itemize}

\begin{tcolorbox}[breakable,colback=teal!4,colframe=teal!35!black,title={Before you move on}]
What does \sk{पंकः} mean? (\textquotedblleft{}Mud\textquotedblright{}.) Say it once more.
\end{tcolorbox}
